% Chapter 02.4: তারযুক্ত মাধ্যম
\begingroup
\chTwoReadableTopicStyle
\section{তারযুক্ত মাধ্যম}

\subsection{Lab থেকে Backbone: তিন ধরনের Cable}
কলেজের কম্পিউটার থেকে সুইচ পর্যন্ত কম খরচে টুইস্টেড পেয়ার ব্যবহার করা যায়। কেবল টিভির বিতরণে কোঅ্যাক্সিয়াল এবং ভবনের মধ্যে বা দূরের প্রধান সংযোগে অপটিক্যাল ফাইবার ব্যবহৃত হয়। তারযুক্ত মাধ্যমের শ্রেণিবিভাগ আগের অংশে আছে। এখানে প্রতিটি কেবলের গঠন, কাজ ও সীমাবদ্ধতা ব্যাখ্যা করা হয়েছে।
\chTwoPair{\chTwoImage{chapters/Chapter02/images/wired_media_overview.png}{কেবলের গঠন ও আলোকপথ}}{\chTwoImage{chapters/Chapter02/images/wired_media_types_tree.png}{তারযুক্ত মাধ্যমের তিনটি পরিচিত ধরন}}
\subsection{Twisted Pair Cable}

\begin{definitionbox}{Twisted Pair Cable}
দুটি অন্তরক আবরণযুক্ত তামার তার পরস্পর পেঁচিয়ে তৈরি ডেটা যোগাযোগের কেবলকে টুইস্টেড পেয়ার কেবল বলে।
\end{definitionbox}

কম্পিউটার নেটওয়ার্কে টুইস্টেড পেয়ার বহুল ব্যবহৃত তামার কেবল। এতে তামার তারগুলো জোড়ায় জোড়ায় পেঁচানো থাকে। পেঁচানো বিন্যাস বিদ্যুৎচৌম্বক বিঘ্ন ও পাশের জোড়ার সংকেতের প্রভাব বা ক্রসটক কমায়। অফিসের LAN, বিদ্যালয়ের কম্পিউটার ল্যাব, রাউটার থেকে কম্পিউটারের সংযোগ এবং টেলিফোন লাইনে এই কেবল ব্যবহৃত হয়।


\chTwoPair{\begin{keypointbox}{Twisted pair cable-এর বৈশিষ্ট্য}
\begin{itemize}
\item তামার তার দিয়ে তৈরি।
\item তারগুলো জোড়ায় জোড়ায় পেঁচানো থাকে।
\item স্থাপন সহজ এবং খরচ কম।
\item LAN-এ সবচেয়ে বেশি ব্যবহৃত কেবল।
\item সাধারণত RJ45 কানেক্টরের সঙ্গে ব্যবহৃত হয়।
\item স্বল্প থেকে মাঝারি দূরত্বের যোগাযোগে উপযোগী।
\item বিদ্যুৎচৌম্বক বিঘ্নে প্রভাবিত হতে পারে।
\end{itemize}
\end{keypointbox}}{\chTwoImage{chapters/Chapter02/images/twisted_pair_cable_structure.png}{টুইস্টেড পেয়ার কেবলের গঠন}}


\begin{examplebox}{বাস্তব উদাহরণ}
বাড়ির Wi-Fi রাউটার থেকে ডেস্কটপ কম্পিউটারে ব্যবহৃত Ethernet কেবল সাধারণত টুইস্টেড পেয়ার। কম্পিউটার ল্যাবে কম্পিউটার থেকে সুইচের সংযোগেও এই কেবল ব্যবহৃত হয়।
\end{examplebox}

\subsection{Twisted Pair Cable কেন পেঁচানো থাকে?}

ক্রসটক হলো একটি তারের জোড়া বা চ্যানেলের সংকেত অনাকাঙ্ক্ষিতভাবে পাশের জোড়া বা চ্যানেলে যুক্ত হওয়া। পেঁচানো দুটি তামার তারে ভারসাম্যযুক্ত সংকেত বহন করা হয়। প্রাপক দুই তারের সংকেতের পার্থক্য পড়ে। ফলে দুই তারে প্রায় সমানভাবে আসা নয়েজের প্রভাব অনেকটা বাদ যায়। তার পেঁচানোর বিন্যাস ক্রসটকও কমায়। EMI-এর মূল ধারণা সংকেতের বিঘ্ন অংশে (অনুচ্ছেদ~\ref{ch2:c01-07}) আছে। শিল্পকারখানার যন্ত্রের পাশে কেবল বসালে ধাতব সুরক্ষা আবরণ ও সঠিক গ্রাউন্ডিংয়ের প্রয়োজন আলাদাভাবে বিচার করতে হয়।

\begin{keypointbox}{পেঁচানোর সুবিধা}
\begin{itemize}
\item ক্রসটক কমায়।
\item বিদ্যুৎচৌম্বক বিঘ্ন কমাতে সাহায্য করে।
\item সংকেতের বিকৃতি কম হয়।
\item কেবল নমনীয় এবং সহজে স্থাপন করা যায়।
\end{itemize}
\end{keypointbox}

\subsection{UTP ও STP Cable}

Twisted pair cable দুই ধরনের:

\begin{itemize}
\item \textbf{UTP}: Unshielded Twisted Pair
\item \textbf{STP}: Shielded Twisted Pair
\end{itemize}

\begin{definitionbox}{UTP Cable}
যে টুইস্টেড পেয়ার কেবলে তারের জোড়ার চারপাশে অতিরিক্ত ধাতব সুরক্ষা আবরণ থাকে না, তাকে UTP বা Unshielded Twisted Pair কেবল বলে।
\end{definitionbox}

\begin{definitionbox}{STP Cable}
বিদ্যুৎচৌম্বক বিঘ্ন কমানোর জন্য অতিরিক্ত ধাতব সুরক্ষা আবরণযুক্ত টুইস্টেড পেয়ার কেবলকে STP বা Shielded Twisted Pair কেবল বলে।
\end{definitionbox}

\begingroup
\begin{longtable}{|p{\dimexpr 0.22\linewidth-2\tabcolsep-\arrayrulewidth\relax}|p{\dimexpr 0.30\linewidth-2\tabcolsep-\arrayrulewidth\relax}|p{\dimexpr 0.30\linewidth-2\tabcolsep-\arrayrulewidth\relax}|}
\hline
\textbf{বিষয়} & \textbf{UTP} & \textbf{STP} \\
\hline
\endfirsthead
\hline
\textbf{বিষয়} & \textbf{UTP} & \textbf{STP} \\
\hline
\endhead

পূর্ণরূপ & Unshielded Twisted Pair & Shielded Twisted Pair \\
\hline
সুরক্ষা আবরণ & বাড়তি ধাতব আবরণ নেই & ধাতব আবরণ আছে \\
\hline
খরচ & কম & তুলনামূলক বেশি \\
\hline
স্থাপন & সহজ & কিছুটা জটিল \\
\hline
বিঘ্ন থেকে সুরক্ষা & তুলনামূলক কম & সঠিক স্থাপনে তুলনামূলক বেশি \\
\hline
ব্যবহার & সাধারণ LAN, অফিস ও বাসার নেটওয়ার্ক & শিল্পকারখানা ও বিঘ্নপ্রবণ এলাকা \\
\hline
\end{longtable}
\endgroup

\begin{mistakebox}{UTP বনাম STP}
UTP ও STP একই কেবল পরিবারের দুটি ধরন। উভয়ই টুইস্টেড পেয়ার; পার্থক্য হলো অতিরিক্ত ধাতব সুরক্ষা আবরণ আছে কি না।
\end{mistakebox}

টুইস্টেড পেয়ার নমনীয় এবং RJ45 দিয়ে সহজে সংযুক্ত করা যায়। বাইরে থেকে কেবল ক্ষতিগ্রস্ত হলে সংযোগ নষ্ট হয়। গুরুত্বপূর্ণ ডেটা রক্ষায় ব্যবহারকারীর প্রবেশাধিকার নিয়ন্ত্রণও প্রয়োজন।

\subsection{Coaxial Cable}

\begin{definitionbox}{Coaxial Cable}
কেন্দ্রীয় পরিবাহী, অন্তরক স্তর, ধাতব সুরক্ষা আবরণ ও বাইরের আবরণ নিয়ে গঠিত কেবলকে কোঅ্যাক্সিয়াল কেবল বলে।
\end{definitionbox}

কোঅ্যাক্সিয়াল কেবলের কেন্দ্রীয় পরিবাহী ও বাইরের ধাতব সুরক্ষা আবরণ একই অক্ষ বরাবর থাকে। এই গঠন থেকেই নামটি এসেছে। কেন্দ্রীয় তামার পরিবাহী সংকেত বহন করে। তার চারপাশে অন্তরক এবং অন্তরকের বাইরে ধাতব সুরক্ষা আবরণ থাকে। সুরক্ষা আবরণ বাইরের নয়েজ ও বিদ্যুৎচৌম্বক বিঘ্ন কমায়। কেবল টিভি, ব্রডব্যান্ড কেবল বিতরণ এবং পুরোনো Ethernet নেটওয়ার্কে কোঅ্যাক্সিয়াল ব্যবহৃত হয়েছে।


\chTwoPair{\begin{keypointbox}{Coaxial cable-এর গঠন}
\begin{itemize}
\item \textbf{Central conductor}: সংকেত বহন করে।
\item \textbf{Insulator}: পরিবাহীকে ধাতব সুরক্ষা আবরণ থেকে আলাদা রাখে।
\item \textbf{Metallic shield}: বাইরের নয়েজ ও বিঘ্ন কমায়।
\item \textbf{Outer jacket}: কেবলকে বাইরের ক্ষতি থেকে রক্ষা করে।
\end{itemize}
\end{keypointbox}}{\chTwoImage{chapters/Chapter02/images/coaxial_cable_structure.png}{কোঅ্যাক্সিয়াল কেবলের গঠন}}


\begin{keypointbox}{Coaxial cable-এর বৈশিষ্ট্য}
\begin{itemize}
\item তামার পরিবাহী ব্যবহার করে।
\item ধাতব সুরক্ষা আবরণ থাকায় টুইস্টেড পেয়ারের তুলনায় নয়েজ কম হতে পারে।
\item কেবল টিভি ও ব্রডব্যান্ড বিতরণে ব্যবহৃত হয়।
\item কিছু ক্ষেত্রে টুইস্টেড পেয়ারের তুলনায় বেশি ব্যান্ডউইথ দিতে পারে।
\item টুইস্টেড পেয়ারের তুলনায় স্থাপন কিছুটা কঠিন।
\item আধুনিক LAN-এ টুইস্টেড পেয়ার ও অপটিক্যাল ফাইবার বেশি ব্যবহৃত হওয়ায় কোঅ্যাক্সিয়ালের ব্যবহার কমেছে।
\end{itemize}
\end{keypointbox}

\begin{examplebox}{বাস্তব উদাহরণ}
কেবল টিভির সংযোগ, পুরোনো ব্রডব্যান্ড কেবল লাইন এবং পুরোনো বাস টপোলজির Ethernet নেটওয়ার্কে কোঅ্যাক্সিয়াল কেবল ব্যবহৃত হতো।
\end{examplebox}

কোঅ্যাক্সিয়াল কেবল টেকসই, তবে তুলনামূলক মোটা ও কম নমনীয়। ফাইবারের মতো উচ্চ বহনক্ষমতা প্রয়োজন হলে শুধু কোঅ্যাক্সিয়ালের ধাতব সুরক্ষা আবরণ দিয়ে সেই চাহিদা মেটে না।

\subsection{Optical Fiber Cable}

\begin{definitionbox}{Optical Fiber Cable}
কাচ বা প্লাস্টিকের তন্তু দিয়ে আলোক সংকেতের সাহায্যে ডেটা পাঠানোর কেবলকে অপটিক্যাল ফাইবার কেবল বলে।
\end{definitionbox}

দূরের প্রধান সংযোগ, FTTH এবং সমুদ্রতলের কেবলে অপটিক্যাল ফাইবার গুরুত্বপূর্ণ। বাইরের EMI ফাইবারের আলোকপথে সরাসরি প্রভাব ফেলে না। কোর ও ক্ল্যাডিংয়ের গঠন আলোকে নির্দিষ্ট পথে পরিচালনা করে।


\chTwoPair{\begin{keypointbox}{Optical fiber cable-এর গঠন}
\begin{itemize}
\item \textbf{Core}: আলোক সংকেত চলাচলের প্রধান পথ।
\item \textbf{Cladding}: কোরের চারপাশে থাকে এবং আলোকে কোরের মধ্যে রাখতে সাহায্য করে।
\item \textbf{Coating}: ফাইবারকে চাপ, আঘাত ও অন্যান্য যান্ত্রিক ক্ষতি থেকে রক্ষা করে।
\item \textbf{Outer jacket}: কেবলের বাইরের সুরক্ষা স্তর।
\end{itemize}
\end{keypointbox}}{\chTwoImage{chapters/Chapter02/images/optical_fiber_cable_structure.png}{অপটিক্যাল ফাইবার কেবলের গঠন}}


\subsection{Optical Fiber Cable কীভাবে কাজ করে?}

প্রেরকের LED বা লেজার বৈদ্যুতিক ডেটাকে আলোক স্পন্দনে রূপান্তর করে। ফাইবার ওই আলো বহন করে। গ্রাহকের ফটোডিটেক্টর আলো থেকে বৈদ্যুতিক সংকেত উদ্ধার করে।
\textbf{পূর্ণ অভ্যন্তরীণ প্রতিফলন:} কোরের প্রতিসরণাঙ্ক ক্ল্যাডিংয়ের চেয়ে বেশি। কোর থেকে সীমানায় আপতিত আলোর কোণ অভিলম্বের সঙ্গে মাপা হয়। এই কোণ সংকট কোণের চেয়ে বেশি হলে আলো কোরের মধ্যে প্রতিফলিত হয়। ফাইবারের আলোকপথ বোঝাতে এই সহজ রশ্মি মডেল ব্যবহৃত হয়। মোড হলো আলো বহনের অনুমোদিত ধরন। কোর ছোট হলেই আলো কেবল একটি সোজা সরু রেখায় চলে—মোডের অর্থ এমন নয়।

\chTwoPair{\textbf{Multicomponent Glass কেন?}\par
একাধিক উপাদান মিশিয়ে তৈরি কাচকে Multicomponent Glass বলে। উপাদানের অনুপাত বদলে কাচের প্রতিসরণাঙ্ক নিয়ন্ত্রণ করা যায়। যেমন, সিলিকায় জার্মানিয়া যোগ করলে কোরের প্রতিসরণাঙ্ক বাড়ানো যায়। কোর ও ক্ল্যাডিংয়ের উপযুক্ত প্রতিসরণাঙ্কের পার্থক্য আলোকে পরিচালনা করে। Multicomponent বলতে একটি কেবলে অনেক ফাইবার বা একটি ফাইবারে অনেক মোড বোঝায় না।}{\textbf{তীক্ষ্ণ বাঁকে সমস্যা কেন?}\par
খুব ছোট ব্যাসার্ধে ফাইবার বাঁকালে কোরে আলো আবদ্ধ থাকার শর্ত বদলায়। কিছু আলো বাইরে বেরিয়ে যাওয়ায় প্রাপকের কাছে কম আলো পৌঁছায়। অতিরিক্ত চাপ বা ভাঙনেও ক্ষতি হয়। তাই কেবল বাঁকানোর সময় নির্ধারিত ন্যূনতম ব্যাসার্ধ বা Minimum Bend Radius মানতে হয়। কেবল U-আকৃতি হলেই সংযোগ বন্ধ হয় না; বাঁকের ব্যাসার্ধ ও ফাইবারের নকশা গুরুত্বপূর্ণ। Bend-insensitive Fiber-এ বাঁকের কারণে ক্ষতি কমানো যায়।}

\chTwoPair{\begin{keypointbox}{Optical fiber cable-এর বৈশিষ্ট্য}
\begin{itemize}
\item আলোক সংকেত দিয়ে ডেটা পাঠায়।
\item ব্যান্ডউইথ খুব বেশি।
\item দীর্ঘ দূরত্বের যোগাযোগে উপযোগী।
\item বিদ্যুৎচৌম্বক বিঘ্নে প্রভাবিত হয় না।
\item সংকেত ক্ষয় তুলনামূলক কম।
\item ডেটার নিরাপত্তা তুলনামূলক বেশি।
\item স্থাপন ও রক্ষণাবেক্ষণের খরচ বেশি।
\item কেবল বাঁকা বা ভেঙে গেলে সমস্যা হতে পারে।
\end{itemize}
\end{keypointbox}}{\chTwoImage{chapters/Chapter02/images/optical_fiber_total_internal_reflection.png}{অপটিক্যাল ফাইবারে পূর্ণ অভ্যন্তরীণ প্রতিফলন}}


\begin{examplebox}{বাস্তব উদাহরণ}
FTTH ব্রডব্যান্ড, ইন্টারনেটের প্রধান সংযোগ, সমুদ্রতলের কেবল, ডেটা সেন্টারের সংযোগ এবং মোবাইল নেটওয়ার্কের প্রধান সংযোগে অপটিক্যাল ফাইবার ব্যবহৃত হয়।
\end{examplebox}


\subsection{Single-mode ও Multimode: Core, Mode ও দূরত্ব}
একটি অনুমোদিত আলোক মোড বহনকারী ফাইবারকে Single-mode এবং একাধিক আলোক মোড বহনকারী ফাইবারকে Multimode বলে। নিচের ছকে কোরের মাপ, আলোর উৎস ও উপযোগিতা তুলনা করা হয়েছে।
\begin{tabularx}{\linewidth}{|p{.22\linewidth}|X|X|}
\hline
\textbf{বিষয়} & \textbf{Single-mode} & \textbf{Multimode}\\\hline
Core & ছোট; প্রচলিত প্রায় ৯ $\mu$m; পরীক্ষায় প্রচলিত মান ৮–১২ $\mu$m Range & বড়; প্রচলিত ৫০ বা ৬২.৫ $\mu$m; পরীক্ষায় প্রচলিত মান ৫০–১০০ $\mu$m Range\\\hline
মোড & একটি পরিচালিত আলোক মোড & একাধিক পরিচালিত আলোক মোড\\\hline
আলোর উৎস & সাধারণত লেজার & LED বা লেজার; আধুনিক দ্রুত সংযোগে লেজারও ব্যবহৃত হয়\\\hline
উপযোগিতা & দীর্ঘ দূরত্ব ও ব্যাকবোন & তুলনামূলক স্বল্প দূরত্বে LAN বা ডেটা সেন্টার\\\hline
সীমা & ট্রান্সসিভার ও সংযোগে দক্ষতা দরকার; যন্ত্রের খরচ বেশি হতে পারে & বিভিন্ন মোডের বিলম্বের পার্থক্যে পালস ছড়ায়; কাছের সংযোগ স্থাপন তুলনামূলক সহজ\\\hline
\end{tabularx}
ফাইবার হালকা ও সরু। তবে মেরামত, কানেক্টর বসানো ও ফাইবারের প্রান্ত জোড়া লাগানো বা Splicing-এর জন্য দক্ষতা লাগে। উপযুক্ত যন্ত্রের খরচও বিবেচনা করতে হয়।

\subsection{Twisted Pair, Coaxial ও Optical Fiber-এর তুলনা}

\begingroup
\begin{longtable}{|p{\dimexpr 0.18\linewidth-2\tabcolsep-\arrayrulewidth\relax}|p{\dimexpr 0.23\linewidth-2\tabcolsep-\arrayrulewidth\relax}|p{\dimexpr 0.23\linewidth-2\tabcolsep-\arrayrulewidth\relax}|p{\dimexpr 0.24\linewidth-2\tabcolsep-\arrayrulewidth\relax}|}
\hline
\textbf{বিষয়} & \textbf{Twisted Pair} & \textbf{Coaxial Cable} & \textbf{Optical Fiber} \\
\hline
\endfirsthead
\hline
\textbf{বিষয়} & \textbf{Twisted Pair} & \textbf{Coaxial Cable} & \textbf{Optical Fiber} \\
\hline
\endhead

সংকেতের ধরন & বৈদ্যুতিক সংকেত & বৈদ্যুতিক সংকেত & আলোক সংকেত \\
\hline
উপাদান & তামার তার & কেন্দ্রীয় তামার পরিবাহী & কাচ বা প্লাস্টিকের ফাইবার \\
\hline
ব্যান্ডউইথ & কেবলের শ্রেণি অনুযায়ী কম থেকে মাঝারি বা উচ্চ & প্রচলিত তুলনায় মাঝারি & উচ্চ বহনক্ষমতা \\
\hline
দূরত্ব & স্বল্প বা মাঝারি & মাঝারি & দীর্ঘ দূরত্ব \\
\hline
বিঘ্ন & বেশি প্রভাবিত হতে পারে & সুরক্ষা আবরণের কারণে কম & আলোকপথ EMI দ্বারা প্রভাবিত নয় \\
\hline
খরচ & কম & মাঝারি & তুলনামূলক বেশি \\
\hline
স্থাপন & সহজ & মাঝারি জটিলতা & তুলনামূলক কঠিন \\
\hline
সাধারণ ব্যবহার & LAN, অফিস ও বাসার নেটওয়ার্ক & কেবল টিভি, পুরোনো ব্রডব্যান্ড & FTTH, ব্যাকবোন, ডেটা সেন্টার \\
\hline
কানেক্টরের পরিচয় & RJ45 & BNC/F-type & ফাইবার কানেক্টর \\
\hline
\end{longtable}
\endgroup

\subsection{তারযুক্ত মাধ্যম নির্বাচন করার কৌশল}

সব পরিস্থিতিতে একই কেবল উপযুক্ত নয়। নেটওয়ার্কের নকশা করার সময় কাজের প্রয়োজন অনুযায়ী কেবল নির্বাচন করতে হয়।

\begingroup
\begin{longtable}{|p{\dimexpr 0.30\linewidth-2\tabcolsep-\arrayrulewidth\relax}|p{\dimexpr 0.26\linewidth-2\tabcolsep-\arrayrulewidth\relax}|p{\dimexpr 0.30\linewidth-2\tabcolsep-\arrayrulewidth\relax}|}
\hline
\textbf{পরিস্থিতি} & \textbf{উপযুক্ত Cable} & \textbf{কারণ} \\
\hline
\endfirsthead
\hline
\textbf{পরিস্থিতি} & \textbf{উপযুক্ত Cable} & \textbf{কারণ} \\
\hline
\endhead

বিদ্যালয়ের ল্যাব & Twisted pair cable & কম খরচ, সহজ স্থাপন, LAN-এর উপযোগী \\
\hline
কেবল টিভি বিতরণ & Coaxial cable & ভিডিও সংকেত বিতরণ ও সুরক্ষা আবরণের সুবিধা \\
\hline
উচ্চগতির ব্রডব্যান্ড ব্যাকবোন & Optical fiber cable & বেশি ব্যান্ডউইথ, দীর্ঘ দূরত্ব, কম বিঘ্ন \\
\hline
নয়েজযুক্ত শিল্প এলাকা & STP cable বা fiber optic & বিঘ্ন থেকে সুরক্ষা বেশি \\
\hline
ডেটা সেন্টারগুলোর সংযোগ & Optical fiber cable & উচ্চগতি ও নির্ভরযোগ্য যোগাযোগ \\
\hline
বাসার রাউটার থেকে কম্পিউটার & UTP cable & কম খরচে সহজ Ethernet সংযোগ \\
\hline
\end{longtable}
\endgroup

\subsection{Cable-এর সংখ্যা, Connector ও দূরত্ব}
\label{ch2:wired-numeric}
\begin{multicols}{2}
\textbf{Ethernet UTP Cable-এ কয়টি তার থাকে?} প্রচলিত চার-পেয়ার UTP-তে ৪ জোড়া বা ৮টি অন্তরক আবরণযুক্ত তামার তার থাকে। জোড়াগুলোর রং সাধারণত কমলা, সবুজ, নীল ও বাদামি। প্রতিটি রঙের তারের সঙ্গে সাদা ডোরাযুক্ত সঙ্গী থাকে। পেঁচানো জোড়া বাইরের নয়েজ ও পারস্পরিক ক্রসটক কমায়।\par
\textbf{RJ45 ও RJ11 কি একই Connector?} Ethernet-এ প্রচলিত ৮টি সংযোগবিন্দুর কানেক্টরকে সাধারণত RJ45 বলা হয়। এর গঠন 8P8C। টেলিফোন লাইনে প্রচলিত ছোট কানেক্টর হলো RJ11। কানেক্টরের নাম এক হলেই সব কেবলের গতি সমান হয় না।\par
\end{multicols}
\textbf{UTP ও STP:} UTP-তে অতিরিক্ত ধাতব সুরক্ষা আবরণ নেই। STP-তে তারের জোড়া বা পুরো কেবলে এই আবরণ থাকে। সুরক্ষা আবরণ ও সঠিক গ্রাউন্ডিং EMI কমায়, তবে খরচ ও স্থাপনের জটিলতা বাড়ায়। প্রচলিত টুইস্টেড-পেয়ার Ethernet চ্যানেল প্রায় ১০০ m পর্যন্ত নকশা করা হয়। কেবলের শ্রেণি, মান, প্রান্তের সংযোগ এবং যন্ত্রের ডেটা হার মিলতে হবে। Cat 3-এর ঐতিহাসিক ব্যবহার ১০ Mbps, Cat 5-এর ১০০ Mbps এবং Cat 5e-এর Gigabit Ethernet পরিচিত উদাহরণ। সব টুইস্টেড পেয়ার কেবলের গতি এক নয়।

\textbf{১০০ m কীভাবে গোনা হয়?} প্রচলিত Structured Cabling-এ প্রায় ৯০ m স্থায়ী সংযোগ ও মোট ১০ m Patch Cord মিলিয়ে ১০০ m চ্যানেল হয়। এই দৈর্ঘ্য একটি তামার সংযোগের; পুরো ক্যাম্পাস নেটওয়ার্কের নয়। সমর্থিত Cat 6A চ্যানেলে ১০ Gbps গতির সংযোগও ১০০ m পর্যন্ত পরিকল্পনা করা যায়। এজন্য কেবল, প্রান্তের সংযোগ এবং দুই প্রান্তের পোর্ট সামঞ্জস্যপূর্ণ হতে হবে।

\subsection{Coaxial-এর Rate, Thinnet ও Thicknet}
\begin{keypointbox}{Rate-এর দুটি আলাদা ঐতিহাসিক প্রসঙ্গ}
নির্দিষ্ট কোঅ্যাক্সিয়াল ব্যবস্থার প্রচলিত উদাহরণে \textbf{২০০ Mbps} হার দেওয়া হয়। নিচের \textbf{10BASE2/10BASE5 Ethernet-এর হার ১০ Mbps}। কেবলের গঠন ও নির্দিষ্ট যোগাযোগব্যবস্থার ডেটা হার আলাদা বিষয়। তাই ২০০ Mbps-কে সব কোঅ্যাক্সিয়াল সংযোগের সর্বোচ্চ ক্ষমতা বলা যায় না।
\end{keypointbox}
\begin{tabularx}{\linewidth}{|p{.22\linewidth}|X|X|}
\hline
\textbf{বিষয়} & \textbf{Thinnet} & \textbf{Thicknet}\\\hline
Ethernet নাম & 10BASE2 & 10BASE5\\\hline
Cable & পাতলা 50-ohm Coaxial & মোটা 50-ohm Coaxial\\\hline
পুরোনো Cable-এর আনুমানিক ব্যাস & ০.২৫ inch; প্রায় ৬ mm & ০.৫ inch; প্রায় ১২ mm\\\hline
Data Rate & ১০ Mbps & ১০ Mbps\\\hline
এক Segment-এর সীমা & ১৮৫ m; নামের 2 প্রায় ২০০ m বোঝায় & ৫০০ m\\\hline
ঐতিহাসিক সংযোগ & BNC ও T-Connector & Tap/Transceiver; মোটা Backbone\\\hline
\end{tabularx}
10 মানে ১০ Mbps এবং BASE মানে Baseband Signaling। কেবল টিভিতে সাধারণত 75-ohm কোঅ্যাক্সিয়াল ও F-type কানেক্টর ব্যবহৃত হয়। এটি 10BASE2 Ethernet-এর BNC সংযোগ থেকে আলাদা।

\subsection{Fiber-এর Mode, Core ও আলোর উৎস}
\textbf{ফাইবার সংযোগের নির্দিষ্ট উদাহরণ:} ১০ Gbps Ethernet-এ 10GBASE-SR মডিউল ও OM3 Multimode Fiber দিয়ে প্রায় ৩০০ m এবং OM4 দিয়ে প্রায় ৪০০ m সংযোগের মান নির্ধারিত আছে। 10GBASE-LR ও উপযুক্ত Single-mode Fiber দিয়ে প্রায় ১০ km সংযোগের মান পাওয়া যায়। এগুলো নির্দিষ্ট মডিউল ও ফাইবারের সীমা; সব ফাইবারের সর্বোচ্চ দূরত্ব এক নয়। ডেটা হার সমান হলেও মোড ও যন্ত্র বদলালে কার্যকর দূরত্ব বদলায়।

\textbf{আলোক যোগাযোগের ইতিহাস:} Charles K. Kao ও George Hockham ১৯৬৬ সালে কম ক্ষয়যুক্ত কাচের ফাইবার দিয়ে দূরের যোগাযোগের সম্ভাবনা ব্যাখ্যা করেন। Kao-কে অপটিক্যাল ফাইবার যোগাযোগের জনক বলা হয়। এই অবদানের জন্য তিনি ২০০৯ সালে পদার্থবিজ্ঞানে নোবেল পুরস্কার পান। গবেষণা, কার্যকর কেবল তৈরি ও বাজারে ব্যবহার পৃথক সময়ে বিকশিত হয়েছে।

\textbf{আলোর উৎসের নাম:} LED = Light Emitting Diode; LASER = Light Amplification by Stimulated Emission of Radiation। আলোর উৎস, ফাইবারের মোড ও গ্রাহক যন্ত্রের সামঞ্জস্যে সংযোগ তৈরি হয়। LED বা লেজার আলোর উৎস; কেবল নয়। প্রেরক ও গ্রাহকের রূপান্তরের উদাহরণ পরের বাস্তব সমস্যায় আছে।
\textbf{আলোর তরঙ্গদৈর্ঘ্য:} পরীক্ষায় প্রচলিত আনুমানিক মান হিসেবে \textbf{প্রায় ১৫০০ nm} পাওয়া যায়। প্রচলিত নিকট-অবলোহিত আলোর ব্যবহৃত পরিসর প্রায় ৮৫০, ১৩১০ ও ১৫৫০ nm। Multimode আলোচনায় ১৩০০ nm-ও ব্যবহৃত হয়। ব্যবস্থাভেদে তরঙ্গদৈর্ঘ্য ও ডেটা হার বদলায়।
\chTwoPair{\textbf{Step-index Fiber:} কোরের মধ্যে প্রতিসরণাঙ্ক প্রায় সমান থাকে। কোর থেকে ক্ল্যাডিংয়ে গেলে প্রতিসরণাঙ্ক হঠাৎ কমে। Multimode Step-index-এ আলোকপথগুলোর দৈর্ঘ্য ভিন্ন। ফলে একই আলোক স্পন্দনের অংশগুলো পৃথক সময়ে পৌঁছে ছড়িয়ে পড়ে।}{\textbf{Graded-index Fiber:} কোরের কেন্দ্র থেকে বাইরের দিকে প্রতিসরণাঙ্ক ধীরে কমে। বাইরের দীর্ঘ পথে আলো তুলনামূলক দ্রুত চলে। এতে Multimode-এর বিভিন্ন পথে বিলম্বের পার্থক্য কমে। Single/Multimode দিয়ে মোডের সংখ্যা এবং Step/Graded দিয়ে প্রতিসরণাঙ্কের বিন্যাস বোঝানো হয়।}
\textbf{অধাতব আলোকপথ:} কাচ বা প্লাস্টিকের ফাইবার বিদ্যুৎ বহনকারী ধাতব তার নয়। আলো দিয়ে ডেটা বহন করায় বাইরের EMI ফাইবারের সংকেতে সরাসরি প্রভাব ফেলে না। তবে সুরক্ষার জন্য কোনো ফাইবার কেবলের বাইরের শক্ত আবরণে ধাতু থাকতে পারে। আলোকপথ অধাতব হওয়া মানে প্রতিটি ফাইবার কেবলের সব অংশ ধাতুমুক্ত নয়।

\textbf{সুবিধার কারণ:} উচ্চ বহনক্ষমতা, কম সংকেত ক্ষয় এবং EMI-মুক্ত আলোকপথ ফাইবারকে প্রধান সংযোগে উপযোগী করে। তবে ফাইবারে সংকেত ক্ষয় শূন্য নয়। সব তামার সংকেতের তুলনায় ফাইবারে আলো ছড়িয়ে পৌঁছানোর গতি সবসময় বেশি—এমনও নয়। তীক্ষ্ণ বাঁক, কানেক্টর ও জোড়ার কারণে ক্ষয় বিবেচনা করতে হয়। যেমন, ২ km দূরের ভবনে সংযোগ দিতে উপযুক্ত ফাইবার মোড, দুই প্রান্তের ট্রান্সসিভার ও জোড়া লাগানোর ব্যবস্থা নির্বাচন করতে হয়। পুরো পথে সংকেতের শক্তি ও ক্ষয়ের হিসাব বা Link Budget-ও মেলাতে হয়। শুধু ফাইবার ব্যবহার করলেই যেকোনো দূরত্বে যেকোনো ডেটা হার নিশ্চিত হয় না।

% Original topic-level problem, solution and MCQ reasoning.
\subsection{বাস্তব সমস্যা: Workshop-এর Noise ও দূরের Server}
\chTwoPair{\textbf{সমস্যা:} কারখানার মোটরের পাশে সংকেতে ত্রুটি হচ্ছে। অন্য ভবনে সার্ভার আছে। সেখানে এবং আরও দূরে ডেটা পাঠাতে হবে।\par
\textbf{সমাধান:} প্রথমে কেবলের পথ, কানেক্টর ও গ্রাউন্ডিং যাচাই করতে হবে। কাছের তামার সংযোগে উপযুক্ত সুরক্ষা আবরণযুক্ত কেবল সাহায্য করতে পারে। দীর্ঘ বা EMI-প্রবণ সংযোগে ফাইবারের আলোকপথ সুবিধাজনক। ট্রান্সসিভার, ফাইবার জোড়া লাগানোর ব্যবস্থা ও খরচ বিবেচনা করতে হবে।}{
\textbf{প্রেরক ও গ্রাহকের রূপান্তর:} গ্রাহকের ফটোডিটেক্টর, যেমন ফটোডায়োড, আলোকে বৈদ্যুতিক সংকেত বা প্রবাহে রূপান্তর করে। প্রেরকের LED বা লেজার বৈদ্যুতিক ইনপুট থেকে আলো তৈরি করে। দুটি রূপান্তর বিপরীত দিকে ঘটে।\par}
\textbf{একই নাম, আলাদা মাপ:} কোর, ক্ল্যাডিং এবং সুরক্ষা আবরণ বা জ্যাকেটের ব্যাস পৃথক। প্রচলিত কাচের ফাইবারে ক্ল্যাডিংয়ের ব্যাস প্রায় ১২৫ $\mu$m। সুরক্ষা আবরণের ব্যাস আরও বড়। কোরের ৮–১২ বা ৫০–১০০ $\mu$m-এর উদাহরণ পুরো কেবলের ব্যাস নির্দেশ করে না।


\subsection{সাধারণ ভুল ও সংশোধন}

\smallhead{১. Twisted pair cable মানেই শুধু telephone cable মনে করা}
টুইস্টেড পেয়ার কেবল শুধু টেলিফোন লাইনে ব্যবহৃত হয় না। আধুনিক Ethernet LAN-এও এই কেবলের ব্যাপক ব্যবহার রয়েছে।

\smallhead{২. Coaxial ও fiber optic এক মনে করা}
কোঅ্যাক্সিয়াল কেবল বৈদ্যুতিক সংকেত এবং অপটিক্যাল ফাইবার আলোক সংকেত বহন করে। তাই দুটি কেবল এক নয়।

\smallhead{৩. Fiber দ্রুত হওয়ার কারণ শুধু পাতলা হওয়া মনে করা}
ফাইবার শুধু পাতলা হওয়ার কারণে দ্রুত ডেটা পাঠায় না। আলোক সংকেত, উচ্চ ব্যান্ডউইথ, কম বিঘ্ন এবং পূর্ণ অভ্যন্তরীণ প্রতিফলনের সাহায্যে নির্দিষ্ট পথে আলো বহন করা এর কার্যকারিতার মূল কারণ।

\smallhead{৪. UTP সব পরিবেশে best মনে করা}
UTP কম খরচের কেবল হলেও সব পরিবেশে উপযুক্ত নয়। বিদ্যুৎচৌম্বক নয়েজ বেশি হলে STP বা অপটিক্যাল ফাইবার ভালো হতে পারে।

\subsection{MCQ অনুশীলন}

\par
\textbf{MCQ}
\begin{multicols}{2}\begin{enumerate}
\item\chTwoMCQBlock{নিচের কোনটি তারযুক্ত মাধ্যম?
\par ক. রেডিও তরঙ্গ \quad খ. মাইক্রোওয়েভ \par গ. টুইস্টেড পেয়ার কেবল \quad ঘ. অবলোহিত রশ্মি }
\item\chTwoMCQBlock{টুইস্টেড পেয়ার কেবলের পরিবাহী তার সাধারণত কী দিয়ে তৈরি?
\par ক. তামা \quad খ. বায়ু \par গ. চৌম্বক পদার্থের ডিস্ক \quad ঘ. উপগ্রহের প্লেট }
\item\chTwoMCQBlock{UTP-এর পূর্ণরূপ কী?
\par ক. Universal Twisted Pair \quad খ. Unshielded Twisted Pair \par গ. Unified Transfer Protocol \quad ঘ. Ultra Transmission Path }
\item\chTwoMCQBlock{কোঅ্যাক্সিয়াল কেবলের কেন্দ্রীয় পরিবাহী কী বহন করে?
\par ক. ডেটাবাহী বৈদ্যুতিক সংকেত \quad খ. কাগজ \par গ. শুধু শব্দতরঙ্গ \quad ঘ. শুধু তাপ }
\item\chTwoMCQBlock{অপটিক্যাল ফাইবার কোন সংকেত ব্যবহার করে?
\par ক. আলোকসংকেত \quad খ. শব্দসংকেত \par গ. যান্ত্রিক সংকেত \quad ঘ. জলের তরঙ্গ }
\item\chTwoMCQBlock{ফাইবারের আলোকপথে বিদ্যুৎচৌম্বক বিঘ্নের প্রভাব কেমন?
\par ক. খুব বেশি প্রভাবিত হয় \quad খ. প্রভাবিত হয় না \par গ. সবসময় নষ্ট হয় \quad ঘ. শুধু রাতে কাজ করে }
\item\chTwoMCQBlock{UTP-এর তুলনায় STP-এর বাড়তি সুবিধা কোনটি?
\par ক. আবরণ দিয়ে বিদ্যুৎচৌম্বক বিঘ্ন থেকে সুরক্ষা \quad খ. কোনো তার লাগে না \par গ. আলোকসংকেত পাঠায় \quad ঘ. স্যাটেলাইট ব্যবহার করে }
\item\chTwoMCQBlock{উচ্চক্ষমতার ইন্টারনেট ব্যাকবোনের জন্য কোন কেবল সবচেয়ে উপযোগী?
\par ক. অপটিক্যাল ফাইবার \quad খ. শুধু UTP \par গ. টিভি রিমোটের কেবল \quad ঘ. অবলোহিত রশ্মির নল }
\end{enumerate}\end{multicols}
\textbf{উত্তরমালা: ১. গ, ২. ক, ৩. খ, ৪. ক, ৫. ক, ৬. খ, ৭. ক, ৮. ক}
\endgroup
